\scanpage{220}

it merely expresses (in view of (154) and \unsure{(150)} the relation
\[
C=\overbrace{\varepsilon(T+2\nu U)}^{C'}.
\]
Thus the relations (154) (155) express the fact that the two sides of (153) \deleted{\unsure{all} \unsure{their terms}} \inserted{three coefficients} equal \deleted{\illegible{1}} \inserted{\unsure{pairwise}}, \inserted{$\unsure{B=B'},\ C=C',\ D=D'$}, and the equality for the last
\[
A=\overbrace{\varepsilon(R+2\nu S)}^{A'}
\]
then follows from this thanks to the relations \deleted{(149)}, \deleted{which allow us to express}
\[
AD-BC=1,\qquad A'D'-B'C'=1,
\]
which give
\[
A=\frac{1+BC}{D},\qquad A'=\frac{1+B'C'}{D'}
\]
(N.B. we know that $D=D'$ is invertible, by (155)) hence $A=A'$.

We thus find the

\uline{Theorem} For a solution $(\xi,\eta)$ in $\widehat{\Pi}_{0,3}\times\widehat{\Pi}_{0,3}$ of the system of equations (133) (134), again denoting by $\xi,\eta$ elements of $\mathrm{Sl}(2,\widehat{\mathbb Z})$ corresponding to them, the sign ambiguity for $\xi$ being removed by relation (147), we will have the relations (154) and (155), which imply the following consequences:

a) $\mu\inserted{=2\nu+1}$ is invertible, and it is entirely \unsure{determined} by $\xi$ -- it is determined by $\eta$ only up to sign.

\scanpage{221}
\mspage{410}

\uline{Remark}

In view of the relations (149) (150)
\[
AD-BC=1,\qquad \unsure{A|C}+D-B=1,
\]
and of the fact that $D$ is invertible, the matrices $\xi$ are entirely determined by the coefficients $B,D$ --- the pairs $(\xi,\eta)$ of matrices $\in\mathrm{Sl}(2,\widehat{\mathbb Z})$, satisfying
\[
\tag{156}
\left\{\begin{gathered}
\xi\rho(\xi)\rho^2(\xi)=1,\quad
\eta\sigma_\infty(\eta)=1,\quad
\xi=\eta\ell_1^\nu
\qquad(\nu\in\widehat{\mathbb Z}\ \text{\unsure{suitable}})\\
\xi\equiv1\quad(2),\qquad\eta\equiv1\quad(2)
\end{gathered}\right.
\]

\margin{Of course, the introduction of the congruences here is \unsure{entirely} incidental --- the algebraic statement \unsure{is independent} of such conditions...}

are in one-to-one correspondence with the pairs $(B,D)$, $B\in\widehat{\mathbb Z}$, $D\in\widehat{\mathbb Z}^{*}$, \inserted{$B\equiv0\ (2)$}, from which we deduce
\[
\xi=\begin{pmatrix}A&B\\C&D\end{pmatrix},\qquad
\eta=\begin{pmatrix}R&S\\T&U\end{pmatrix},\qquad
\text{and }\nu\in\widehat{\mathbb Z}\text{ by}
\]
\[
\tag{157}
\begin{aligned}
C&=B-D+1,\qquad
A=\frac{1+BC}{D}=\frac1D(1+B^2-BD+B)\\
&&\inserted{{}=\mu(1+B+B^2)-B}
\end{aligned}
\]
\[
\tag{158}
\mu=\frac1D,\qquad\nu=\frac{\mu-1}{2},
\]
\deleted{\unsure{(159)}}
\[
\eta=\xi\ell_1^{-\nu}
=\begin{pmatrix}A&B\\C&D\end{pmatrix}
\begin{pmatrix}1&0\\-2\nu&1\end{pmatrix}
=\begin{pmatrix}A-2\nu B&B\\C-2\nu D&D\end{pmatrix}
\qquad\text{i.e.}
\]

\margin{we have normalized the sign of $\eta$ as \unsure{said} in (156)}

\[
\tag{159}
\text{i.e.}\quad
\unsure{U}=D,\quad S=\unsure{T}=B,\quad
R=A-2\nu B=\mu(1+B^2)=\frac{1+B^2}{D}
\]

Put (compare page~\unsure{398'} (78))
\[
\tag{160}
\underline u_{\xi\deleted{\unsure{,\eta}}}
=\mu\begin{pmatrix}1&B\\0&D\end{pmatrix}
=\begin{pmatrix}\mu&\mu B\\0&1\end{pmatrix}
=\begin{pmatrix}\dfrac1D&B/D\\0&1\end{pmatrix}
\in\mathrm{Gl}(2,\widehat{\mathbb Z})
\]

\margin{\deleted{for the moment we \unsure{have} an element $\xi$ of $\unsure{\widetilde{\widehat{\Pi}}_{0,3}}$, \unsure{which is determined|$\xi$ is determined by} as an element of $\mathrm{Sl}(2,\widehat{\mathbb Z})'$ \unsure{or in|or of} $\mathrm{Sl}(2,\widehat{\mathbb Z})$ \uuline{up to sign}}}

so that knowledge of $\underline u_\xi$ is equivalent to that \inserted{of $B$ and $D$, i.e. to that} of $\xi\in\mathrm{Sl}(2,\widehat{\mathbb Z})$, or again to that of the pairs $(\xi,\eta)$ ($\eta$ normalized by $\eta=+\xi\ell_1^{-\nu}$). (In loc. cit. this element was denoted $u_{\alpha,\beta}$... but it depends only on $\xi\in\widehat{\Pi}_{0,3}$, and even only on $\underline\xi\in\mathrm{Sl}(2,\widehat{\mathbb Z})$ --- and knowledge of it is equivalent to that of $\underline\xi\in\mathrm{Sl}(2,\widehat{\mathbb Z})$ ---)

\margin{\uline{N.B.} We deduce $\xi,\eta$ from knowledge of $\underline u$ by the formulas
\deleted{\illegible{1}}
\[
\tag{161}
\left\{\begin{aligned}
\underline u\cdot{}^t\underline u&=\mu\eta\\
\mu&=\det\underline u
\end{aligned}\right.
\]
whence
\[
\tag{162}
\left\{\begin{aligned}
\eta&=\mu^{-1}\underline u\cdot{}^t\underline u=\mu^{-1}\unsure{,}\\
\xi&=\eta\ell_1^{-\nu}=\mu^{-1}(\underline u\cdot{}^t\underline u)\ell_1^{-\nu}
\end{aligned}\right.
\]
}

\scanpage{222}

\deleted{\unsure{the map}} We will have
\deleted{$\unsure{\mathcal M_{0,3}^{!\sim}(0)}$}
\deleted{\unsure{i.e. $\eta=(\underline u,\sigma_\infty)$}}
\[
(163)\qquad
\underline u_\xi(\sigma_\infty)
=\eta\sigma_\infty
=(\xi \ell_1^{-\nu})\sigma_\infty,
\qquad
\underline u_\xi(\rho)=\xi\rho
\]
\[
(164)\qquad
\underline u_\xi(\varepsilon_0)=\varepsilon_0^\mu
\qquad\text{and}\qquad
u_\xi(\ell_0]=\ell_0^\mu
\]

\margin{The formulas (163) can also be written
\[
\unsure{(\underline u,\sigma_\infty)|[\underline u,\sigma_\infty]}=\eta,
\quad [\underline u,\rho]=\xi
\qquad (163\ \mathrm{bis}),
\]
they give the way to \unsure{reconstruct}
$(\unsure{\xi},\eta)$ using $\underline u_\xi=\underline u$.}

If we consider an arbitrary invertible matrix
\[
\underline u=\begin{pmatrix}L&M\\N&P\end{pmatrix}
\in\mathrm{Gl}(2,\widehat{\mathbb Z})
\]
the relation
\[
\deleted{(\illegible{1})}\qquad
\underline u(\varepsilon_0)=\varepsilon_0^\mu
\quad\text{i.e.}\quad
u\varepsilon_0=\varepsilon_0^\mu
\qquad(\text{where }\mu\in\widehat{\mathbb Z}^{*})
\]
\deleted{\unsure{is equivalent to}}\inserted{is written}
\[
\begin{pmatrix}L&-L+M\\N&-N+P\end{pmatrix}
=\begin{pmatrix}L-\mu N&M-\mu P\\N&P\end{pmatrix}
\]
and is equivalent to: \deleted{\unsure{$N$}}
\[
N=0,\qquad L=\mu P
\]
i.e.
\[
\underline u=\begin{pmatrix}\mu P&M\\0&P\end{pmatrix},
\qquad\text{whence}\quad\det u=\mu P^2
\]

the additional condition
\[
\unsure{(165)}\qquad \det\underline u=\mu
\]
gives the condition
\[
\deleted{(\illegible{1})}\qquad P^2=\underline 1,
\]
the most \unsure{obvious} solution being $P=1$,
which gives \unsure{just} the matrices
\[
\underline u=\begin{pmatrix}\mu&M\\0&1\end{pmatrix}
=\underline u_\xi
\]

\scanpage{223}
\mspage{411}

where $\xi=\begin{pmatrix}A&B\\C&D\end{pmatrix}$
\deleted{\unsure{satisfies}} is \unsure{given|given} in terms of
$B,D$ by $D=\dfrac1\mu$, $B=DM=M/\mu$.

\uline{Theorem} The map
\[
(165)\qquad
\begin{array}{rcl}
\mathcal M_{0,3}^{!\sim}[0]&\longrightarrow&
\mathrm{Gl}(2,\widehat{\mathbb Z})^{\deleted{\unsure{\prime|*}}}\\
\underset{u}{u_{\xi,\eta}}&\longmapsto&\underline u_\xi
\end{array},
\quad\text{or again }u\longmapsto\underline u
\]
$\underline u_\xi$ is given by (160)) is a group homomorphism.

\margin{\uline{N.B.} Since we \unsure{have restricted ourselves to}
$\mathcal M_{0,3}^{!\sim}(0)$, we get a homomorphism into
$\mathrm{Gl}(2,\widehat{\mathbb Z})$,
not just into $\mathrm{Gl}(2,\widehat{\mathbb Z})/\mathord\pm1$.}

\margin{\uline{N.B.} This map factors through the
triangular \unsure{subgroup|subgroup} $\mathcal B$ below (169)}

\uuline{Proof}
Let us consider $u=u_{\xi,\eta}$,
\deleted{$u'=u_{\xi',\eta'}$, we will have \unsure{let}
$u''=u_{\xi'',\eta''}$,}
$\underline u$ is its image in $\mathrm{Gl}(2,\widehat{\mathbb Z})$,
and if $x\in\widehat\Pi_{0,3}$, $\underline x$
\deleted{$\xi''=u'(\xi)\xi'\qquad(\text{in }\widehat\Pi_{0,3})$}
its image in $\mathrm{Sl}(2,\widehat{\mathbb Z})$:
\deleted{we will have \unsure{therefore in}}
\[
(166)\qquad
\underline{u(x)}=\underline u(\underline x)
\left(\stackrel{\text{definition}}{=}
\underline u\,\underline x\,\underline u^{-1}\right)
\]
for it suffices to see it for $x=\sigma_\infty$ and for $x=\rho$,
where this formula reduces \unsure{to} (161). This noted,
if $u=u_{\xi,\eta}\inserted{=u_\xi}$,
$u'=u_{\xi',\eta'}\inserted{=u_{\xi'}}$
are two elements of $\mathcal M_{0,3}^{!\sim}(0)$,
$u''=u_{\xi'',\eta''}=u_{\xi''}$ their product, we will have
\inserted{(139)}
\[
\deleted{\unsure{(167)}}\qquad\xi''=u'(\xi)\xi'
\]
whence
\[
\underline\xi''=\underline{u'(\xi)}\,\underline\xi^{\unsure{\prime}}
=\underline u'(\underline\xi)\,\underline\xi',\quad\text{i.e.}
\]
\[
\deleted{\unsure{(167)}}\quad(167)\qquad
\underline\xi''=\operatorname{int}
(\underline u_{\underline\xi'})(\underline\xi)\mathbin{.}\underline\xi'.
\]
and it suffices to verify that if three solutions
\deleted{\illegible{1}}
\deleted{$\xi,\xi',\xi''$ of the equations}
\[
\xi=\begin{pmatrix}A&B\\C&D\end{pmatrix},
\quad\xi'=\begin{pmatrix}A'&B'\\C'&D'\end{pmatrix},
\quad\xi''=\begin{pmatrix}A''&B''\\C''&D''\end{pmatrix}
\]
\unsure{of the} equations
$AD-BC=1$, $C+D\unsure{-}B=1$, $D\in\widehat{\mathbb Z}^{*}$,
are related by (167), then we have
\[
(168)\qquad
\underline u_{\xi''}=\underline u_{\xi'}\mathbin{.}\underline u_\xi
\]

\scanpage{224}

\inserted{Since $\xi$ and $\underline u=u_\xi$ determine each other
(163bis) it}
\deleted{\illegible{1}} amounts to the same thing to say that if,
on the group
\[
(169)\qquad
\begin{array}{ll}
\mathcal B\subset\mathrm{Gl}(2,\widehat{\mathbb Z})&\\
\Vert&\\[-2pt]
\left\{\begin{pmatrix}\mu&M\\0&1\end{pmatrix}
\ \middle|\ \mu,M\in\widehat{\mathbb Z},\ \mu\text{ invertible}\right\}&
\end{array}
\]
we define (cf. (162))
\[
(170)\qquad
\left\{\begin{aligned}
F_\xi &: \mathcal B\longrightarrow\mathrm{Sl}(2,\widehat{\mathbb Z})\\
F_\xi(\underline u)
&=(\underbrace{\det\underline u}_{\mu})^{-1}
\underbrace{\deleted{\unsure{\underline u\,{}^t\underline u}}}
_{\inserted{[\underline u,\rho]}}
\overset{\deleted{\unsure{\mu-1=2\nu}}}
{\deleted{\unsure{\varepsilon_1^{-(\det\underline u-1)}|\ell_1^{-(\det\underline u-1)}}}}
\end{aligned}\right.
\]
then we have the 1-cocycle relation
\[
(171)\qquad
F_\xi(\underline u'\underline u)
=\operatorname{int}(\underline u')\bigl(F_\xi(\underline u)\bigr)
\mathbin{.}F_\xi(\underline u').
\]
By multiplicativity of $\det$, we are \unsure{reduced}
\deleted{$\unsure{(\underline u,\underline u'\in\mathcal B)}$}
to seeing that
\[
\underline u\longmapsto[\underline u,\rho]
=\underline u\,\rho(\underline u)^{-1}
\]
satisfies the cocycle relation, which is trivial.

\uline{Corollary} The map
\[
(172)\qquad
\begin{array}{ccl}
\mathcal M_{0,3}^{\unsure{!\sim}}
=\mathcal M_{0,3}^{!\sim}[0]\mathbin{.}
\widehat{\mathcal G}_{0,3}^{+}
&\longrightarrow&\mathrm{Gl}(2,\widehat{\mathbb Z})'\\
\deleted{\text{product }\tfrac12\text{ direct}}&&\\
u_\xi\mathbin{.}a&\longrightarrow&(\underline u_\xi)'\mathbin{.}\underline a
\end{array}
\]
where $(\underline u_\xi)'$ \unsure{is the} reduction of $\underline u_\xi$
modulo $\pm1$, and where $\underline a$ denotes the image of $a$ in
$\mathrm{Sl}(2,\widehat{\mathbb Z})'$) is a group homomorphism.

Follows at once from the Theorem and the formulas (163).

\margin{careful, this is \uuline{not} a semidirect product,
since $\mathcal M_{0,3}^{!\sim}(0)\cap
\widehat{\mathcal G}_{0,3}^{+}
=L_0\ne\{1\}$, we must therefore show that the defining formula
does not depend on the way of writing a $u$ in the form
\unsure{$u_\xi a$} --- this amounts to seeing that for
$\lambda\in L_0$, \unsure{these two} images \illegible{2} \unsure{as an element of one or the other}
\unsure{factor}, are the same.}

\scanpage{225}
\mspage{412}

From now on I am going to \unsure{specify|specify} the pair
$(\underline\xi,\underline\eta)$ from the adelic point of view
$(\text{in }\mathrm{Gl}(2,\widehat{\mathbb Z}))$ using
$\underline u_\xi$, which is therefore of the form
\[
(173)\qquad
\underline u=\begin{pmatrix}\mu&M\\0&1\end{pmatrix}
\]
\inserted{with $M\equiv0\ (2)$ if $\underline u$ comes from
$\mathcal M_{0,3}^{!\sim}(0)$}
from which we get $\xi,\eta$ by the formulas \unsure{(163 bis)|(167 bis)},
\[
(174)\qquad
\underline\xi=[\underline u,\underline\rho],\qquad
\underline\eta=[\underline u,\underline\sigma_\infty].
\]

For $\underline u\in\mathcal M_{0,3}^{!\sim}(0)$, defined by
$(\xi,\eta)$, as long as we do not bring in the
possibility of writing $\xi$ and $\eta$ uniquely as
\[
\xi=\alpha\rho(\alpha)^{-1},\quad
\eta=\beta\sigma_\infty(\beta)^{-1}
\]
with $\alpha\in\unsure{\widetilde{\widehat\Pi}_{0,3}|\widehat{\widetilde\Pi}_{0,3}}$,
$\beta\in\widehat\Pi_{0,3}$, the only adelic invariants that we
can associate to it directly in a natural way is the pair
$(\underline\xi,\underline\eta)$, or again \unsure{(what} amounts to the
same thing, but is prettier) $\underline u$.
The elements \deleted{\illegible{1}}
\[
\lambda=\ell_0^n\in L_0\hookrightarrow\mathcal M_{0,3}^{!\sim}[0]
\]
correspond to
\[
\begin{aligned}
(175)\qquad
\xi=\eta=\ell_0^n\ell_1^{-n}
&=\begin{pmatrix}\unsure{1}&-2n\\0&\unsure{1}\end{pmatrix}
\begin{pmatrix}1&0\\-2n&1\end{pmatrix}\\
&=\begin{pmatrix}A=\unsure{1+4n^2}&B=-2n\\C=-2n&D=1\end{pmatrix}
\end{aligned}
\]
\margin{\unsure{please} do the \unsure{forgotten} check
which \unsure{is} referred to in the note \unsure{in the margin}
\unsure{just before}\ldots}
\[
(176)\qquad
\underline\lambda=\unsure{\underline{\underline{\ell_0}}^{\,n}|\underline{\ell_0}^{\,n}}
=\begin{pmatrix}1&-2n\\0&1\end{pmatrix}
\]
moreover we have
\[
(177)\qquad
\underbrace{\begin{pmatrix}1&-2n\\0&1\end{pmatrix}}_{\underline{\ell_0}^{\,n}}
\begin{pmatrix}\mu&M\\0&1\end{pmatrix}
=\begin{pmatrix}\mu&M-2n\\0&1\end{pmatrix}
\]
which shows that \deleted{\illegible{1}}
\deleted{the quotient}\inserted{the quotients of the}
subgroups $\mathcal B'_1$ of $\mathcal B$ (169) consisting of the